
% Takes two (scaled) points in a SCALED coordinate system (x, y and z)
% and outputs the 
\newcommand{\connectPoints}[7]{
    \begingroup
    \pgfmathsetmacro{\pox}{#1}
    \pgfmathsetmacro{\poy}{#2}
    \pgfmathsetmacro{\poz}{#3}

    \pgfmathsetmacro{\ptx}{#4}
    \pgfmathsetmacro{\pty}{#5}
    \pgfmathsetmacro{\ptz}{#6}

    % length between points 1 and 2
    \pgfmathsetmacro{\length}{sqrt((\ptx-\pox)^2+(\pty-\poy)^2+(\ptz-\poz)^2)}

    % tdplot rotated origin
    \coordinate (Shift) at (\pox,\poy,\poz);
    \tdplotsetrotatedcoordsorigin{(Shift)}

    % obtaining "some" unit vectors for our desired plane
    \pgfmathsetmacro\zAng{atan2(\plineny,\plinenx)}
    \pgfmathsetmacro\yAng{atan2(sqrt((\plinenx)^2+(\plineny)^2),(\plinenz))}

    \pgfmathsetmacro{\tdrotx}{\pox + cos(\zAng)*cos(\yAng)}
    \pgfmathsetmacro{\tdroty}{\poy + sin(\zAng)*cos(\yAng)}
    \pgfmathsetmacro{\tdrotz}{\poz - sin(\yAng)}

    
    \pgfmathsetmacro{\zAngTwoNumerator}{
        (\tdrotx - \pox) * (\ptx - \pox) + 
        (\tdroty - \poy) * (\pty - \poy) + 
        (\tdrotz - \poz) * (\ptz - \poz)
    }
    \pgfmathsetmacro{\zAngTwoDenominator}{
        sqrt(
            (\tdrotx - \pox)^2 + 
            (\tdroty - \poy)^2 + 
            (\tdrotz - \poz)^2
        ) * 
        sqrt(
            (\ptx - \pox)^2 + 
            (\pty - \poy)^2 + 
            (\ptz - \poz)^2
        )
    }
    \pgfmathparse{\zAngTwoDenominator==0}
    \ifnum\pgfmathresult=1
    \pgfmathsetmacro{\zAngTwoDenominator}{0.005}
    \fi
    \pgfmathsetmacro{\zAngTwo}{acos(\zAngTwoNumerator/\zAngTwoDenominator)}

    \pgfmathsetmacro{\aa}{\tdrotx - \pox}
    \pgfmathsetmacro{\bb}{\tdroty - \poy}
    \pgfmathsetmacro{\cc}{\tdrotz - \poz}
    \pgfmathsetmacro{\dd}{\ptx - \pox}
    \pgfmathsetmacro{\ee}{\pty - \poy}
    \pgfmathsetmacro{\ff}{\ptz - \poz}
    \tdplotcrossprod(\aa,\bb,\cc)(\dd,\ee,\ff)
    \pgfmathparse{
        \tdplotresx * \CameraX +
        \tdplotresy * \CameraY +
        \tdplotresz * \CameraZ <
        0
    }
    \ifnum\pgfmathresult=1\pgfmathsetmacro{\zAngTwo}{-\zAngTwo}\fi
    


    \tdplotsetrotatedcoords{\zAng}{\yAng}{\zAngTwo}
    \tikzset{recent_rotated_coords/.style = {
        tdplot_rotated_coords,canvas is xy plane at z = 0
    }}
    \begin{scope}[tdplot_rotated_coords,canvas is xy plane at z = 0]
        \pgflowlevelsynccm
        \pvVector[#7]{0,0}{\length,0} 
    \end{scope}
    \endgroup
}

\newcommand{\rotatedDot}[4]{
    \begingroup
    \pgfmathsetmacro{\pox}{\plinescalex*#1}
    \pgfmathsetmacro{\poy}{\plinescaley*#2}
    \pgfmathsetmacro{\poz}{\plinescalez*#3}
    \pgfmathsetmacro{\nx}{\plinescalex*\plinenx}
    \pgfmathsetmacro{\ny}{\plinescaley*\plineny}
    \pgfmathsetmacro{\nz}{\plinescalez*\plinenz}
    \pgfmathsetmacro{\nl}{sqrt((\nx)^2 + (\ny)^2 + (\nz)^2)}
    \pgfmathsetmacro{\nx}{\nx/\nl}
    \pgfmathsetmacro{\ny}{\ny/\nl}
    \pgfmathsetmacro{\nz}{\nz/\nl}
    \pgfmathsetmacro{\tempnx}{1}
    \pgfmathsetmacro{\tempny}{0}
    \pgfmathsetmacro{\tempnz}{-\nx/\nz}
    \tdplotcrossprod(\nx,\ny,\nz)(\tempnx,\tempny,\tempnz)
    \pgfmathsetmacro{\ox}{\tdplotresx}
    \pgfmathsetmacro{\oy}{\tdplotresy}
    \pgfmathsetmacro{\oz}{\tdplotresz}
    \pgfmathsetmacro{\ol}{sqrt((\ox)^2 + (\oy)^2 + (\oz)^2)}
    \pgfmathsetmacro{\ox}{\ox/\ol}
    \pgfmathsetmacro{\oy}{\oy/\ol}
    \pgfmathsetmacro{\oz}{\oz/\ol}
    \tdplotcrossprod(\nx,\ny,\nz)(\ox,\oy,\oz)
    \pgfmathsetmacro{\tx}{\tdplotresx}
    \pgfmathsetmacro{\ty}{\tdplotresy}
    \pgfmathsetmacro{\tz}{\tdplotresz}
    \pgfmathsetmacro{\tl}{sqrt((\tx)^2 + (\ty)^2 + (\tz)^2)}
    \pgfmathsetmacro{\tx}{\tx/\tl}
    \pgfmathsetmacro{\ty}{\ty/\tl}
    \pgfmathsetmacro{\tz}{\tz/\tl}
    % Convert (nx, ny, nz) to spherical angles for tdplot
    \pgfmathsetmacro{\theta}{atan2(\ny,\nx)} % azimuth
    \pgfmathsetmacro{\phi}{acos(\nz)}        % inclination
    % Set rotated coordinate system
    \tdplotsetrotatedcoords{\theta}{\phi}{0}
    \coordinate (Shift) at (\pox,\poy,\poz);
    \tdplotsetrotatedcoordsorigin{(Shift)}
    % Draw in 2D inside 3D plane
    \begin{scope}[tdplot_rotated_coords, canvas is xy plane at z=0]
        \pgflowlevelsynccm
        \pvDot[#4]{0,0}
    \end{scope}
    % \begin{scope}[
    %     ,shift = {(\pox,\poy,\poz)}
    %     ,x = {(\ox,\oy,\oz)}
    %     ,y = {(\tx,\ty,\tz)}
    %     ,z = {(\nx,\ny,\nz)}
    %     %,shift = {(\pox,\poy,\poz)}
    %     ,canvas is xy plane at z = 0
    % ]
    %     \pgflowlevelsynccm
    %     \pvDot[#4]{0,0}; 
    % \end{scope}
    \endgroup
}





%%%%%%%%%%%%%%%%%%%%%%%%%%%%%%%%%%%%%%%%%%%%%%%%%%%%%%%%%%%%%%%%%%%%%%%%%%%%%%%
% We initially declare these functions with meaningless output,
% so they can be redeclared inside a function we will use several
% times, which redefined them locally to produce meaningful output.
\pgfmathdeclarefunction{planeX}{2}{\pgfmathparse{100}}
\pgfmathdeclarefunction{planeY}{2}{\pgfmathparse{100}}
\pgfmathdeclarefunction{planeZ}{2}{\pgfmathparse{100}}
\newcommand{\addintersection}[3]{redefine me}
\newcommand{\getplane}[2][]{
    \begingroup
    \pgfqkeys{/pline}{#1}

    % First we set the scales based on \pline@size.
    % planes can occupy very disproportionate rectangular
    % prisms. So, we scale all of them inso a simple cube of
    % a provided size.
    \pgfmathsetmacro{\plinescalex}{\pline@size/(\pline@xmax-\pline@xmin)}
    \pgfmathsetmacro{\plinescaley}{\pline@size/(\pline@ymax-\pline@ymin)}
    \pgfmathsetmacro{\plinescalez}{\pline@size/(\pline@zmax-\pline@zmin)}
    \global\let\plinescalex\plinescalex
    \global\let\plinescaley\plinescaley
    \global\let\plinescalez\plinescalez
    \global\let\plinesize\pline@size

    % This redefined the plane equations, based on our normal equation.
    \pgfmathredeclarefunction{planeX}{\pgfmathparse{(\pline@nd-\pline@nb*##1-\pline@nc*##2)/\pline@na}}
    \pgfmathredeclarefunction{planeY}{\pgfmathparse{(\pline@nd-\pline@na*##1-\pline@nc*##2)/\pline@nb}}
    \pgfmathredeclarefunction{planeZ}{\pgfmathparse{(\pline@nd-\pline@na*##1-\pline@nb*##2)/\pline@nc}}

    % New segment list with a unique global name
    \NewSegmentList{intersections#2}

    % The add intersection command is dependent on
    \renewcommand{\addintersection}[3]{

        % These need to be made into pgfmathmacros before being appended to the list.
        \pgfmathsetmacro{\intersectx}{##1}
        \pgfmathsetmacro{\intersecty}{##2}
        \pgfmathsetmacro{\intersectz}{##3}
        \pgfmathparse{
            \intersectx>=\pline@xmin && \intersectx<=\pline@xmax &&
            \intersecty>=\pline@ymin && \intersecty<=\pline@ymax &&
            \intersectz>=\pline@zmin && \intersectz<=\pline@zmax
        }
        \ifnum\pgfmathresult=1
            \ExpandArgs{nne}\AddSegment{intersections#2}
                {0}
                {%
                \def\noexpand\pointx{\intersectx}%
                \def\noexpand\pointy{\intersecty}%
                \def\noexpand\pointz{\intersectz}%
                }%
        \fi
    }
    
    % These for loops go through each face
    % and find the points on it's boundaries that the plane intersects.
    \foreach \y in {\pline@ymin,\pline@ymax} {
        \foreach \z in {\pline@zmin,\pline@zmax} {
            \pgfmathsetmacro{\x}{planeX(\y,\z)}
            \addintersection{\x}{\y}{\z}
        }
    }
    \foreach \x in {\pline@xmin,\pline@xmax} {
        \foreach \z in {\pline@zmin,\pline@zmax} {
            \pgfmathsetmacro{\y}{planeY(\x,\z)}
            \addintersection{\x}{\y}{\z}
        }
    }
    \foreach \x in {\pline@xmin,\pline@xmax} {
        \foreach \y in {\pline@ymin,\pline@ymax} {
            \pgfmathsetmacro{\z}{planeZ(\x,\y)}
            \addintersection{\x}{\y}{\z}
        }
    }
    % normalize the normal
    \pgfmathsetmacro{\magnitude}{sqrt((\pline@na)^2 + (\pline@nb)^2 + (\pline@nc)^2)}
    \pgfmathsetmacro{\plinenx}{\pline@na/\magnitude}
    \pgfmathsetmacro{\plineny}{\pline@nb/\magnitude}
    \pgfmathsetmacro{\plinenz}{\pline@nc/\magnitude}
    \global\let\plinenx\plinenx
    \global\let\plineny\plineny
    \global\let\plinenz\plinenz
    \global\let\plinend\pline@nd
    \global\let\plinemagnitude\magnitude

    % obtain non-parallel vector to the normal
    % same as normal, except \nx is shifted by 1
    % I'll just implement it manually

    % use cross product to find perpendicular vector
    % we will normalize it
    \pgfmathsetmacro{\plineux}{crossX(\plinenx,\plineny,\plinenz,\plinenx+1,\plineny,\plinenz)}
    \pgfmathsetmacro{\plineuy}{crossY(\plinenx,\plineny,\plinenz,\plinenx+1,\plineny,\plinenz)}
    \pgfmathsetmacro{\plineuz}{crossZ(\plinenx,\plineny,\plinenz,\plinenx+1,\plineny,\plinenz)}
    \pgfmathsetmacro{\ul}{sqrt((\plineux)^2 + (\plineuy)^2 + (\plineuz)^2)}
    \pgfmathsetmacro{\plineux}{\plineux/\ul}
    \pgfmathsetmacro{\plineuy}{\plineuy/\ul}
    \pgfmathsetmacro{\plineuz}{\plineuz/\ul}
    \global\let\plineux\plineux
    \global\let\plineuy\plineuy
    \global\let\plineuz\plineuz

    % the second planar vector is the cross product
    % of the normal and the first plane vector.
    \pgfmathsetmacro{\plinevx}{crossX(\plinenx,\plineny,\plinenz,\plineux,\plineuy,\plineuz)}
    \pgfmathsetmacro{\plinevy}{crossY(\plinenx,\plineny,\plinenz,\plineux,\plineuy,\plineuz)}
    \pgfmathsetmacro{\plinevz}{crossZ(\plinenx,\plineny,\plinenz,\plineux,\plineuy,\plineuz)}
    \global\let\plinevx\plinevx
    \global\let\plinevy\plinevy
    \global\let\plinevz\plinevz
    \NewSegmentList{sortedintersections#2}
    % obtain the centroid of the slipped plane.
    \pgfmathsetmacro{\centroidx}{0}
    \pgfmathsetmacro{\centroidy}{0}
    \pgfmathsetmacro{\centroidz}{0}
    \pgfmathsetmacro{\numberofpoints}{0}
    \LoopOverSegmentList{intersections#2}
        {
            \pgfmathsetmacro{\centroidx}{\centroidx+\pointx}
            \pgfmathsetmacro{\centroidy}{\centroidy+\pointy}
            \pgfmathsetmacro{\centroidz}{\centroidz+\pointz}
            \pgfmathsetmacro{\numberofpoints}{\numberofpoints+1}
            \global\let\centroidx\centroidx
            \global\let\centroidy\centroidy
            \global\let\centroidz\centroidz
            \global\let\numberofpoints\numberofpoints
        }
    \pgfmathsetmacro{\centroidx}{\centroidx/\numberofpoints}
    \pgfmathsetmacro{\centroidy}{\centroidy/\numberofpoints}
    \pgfmathsetmacro{\centroidz}{\centroidz/\numberofpoints}
    \LoopOverSegmentList{intersections#2}
        {
            \pgfmathsetmacro{\plineax}{
                dotproduct(
                    \pointx-\centroidx
                    ,\pointy-\centroidy
                    ,\pointz-\centroidz
                    ,\plineux,\plineuy,\plineuz
                )
            }
            \pgfmathsetmacro{\plineay}{
                dotproduct(
                    \pointx-\centroidx
                    ,\pointy-\centroidy
                    ,\pointz-\centroidz
                    ,\plinevx,\plinevy,\plinevz
                )
            }
            \pgfmathsetmacro{\anglea}{atan2(\plineay,\plineax)}
            \ExpandArgs{nne}\AddSegment{sortedintersections#2}
                {\anglea}
                {%
                \def\noexpand\pointx{\pointx}%
                \def\noexpand\pointy{\pointy}%
                \def\noexpand\pointz{\pointz}%
                }%
        }
    \SortSegmentList{sortedintersections#2}
    \xdef\plinepathone{}
    \LoopOverSegmentList{sortedintersections#2}
        {
            \xdef\plinepathone{\plinepathone (\plinescalex*\pointx,\plinescaley*\pointy,\plinescalez*\pointz) --}
        }
    \pgfmathparse{\numberofpoints>=3}
    \ifnum\pgfmathresult=1
        \begin{scope}
            \path[spath/save global = plane#2] \plinepathone cycle;
            \path[spath/save global = lower#2] 
                (\plinescalex*\pline@xmin,\plinescaley*\pline@ymin,\plinescalez*\pline@zmin) -- 
                (\plinescalex*\pline@xmax,\plinescaley*\pline@ymin,\plinescalez*\pline@zmin) -- 
                (\plinescalex*\pline@xmax,\plinescaley*\pline@ymax,\plinescalez*\pline@zmin) --
                (\plinescalex*\pline@xmin,\plinescaley*\pline@ymax,\plinescalez*\pline@zmin) -- cycle;
            \path[spath/save global = upper#2] 
                (\plinescalex*\pline@xmin,\plinescaley*\pline@ymin,\plinescalez*\pline@zmax) -- 
                (\plinescalex*\pline@xmax,\plinescaley*\pline@ymin,\plinescalez*\pline@zmax) -- 
                (\plinescalex*\pline@xmax,\plinescaley*\pline@ymax,\plinescalez*\pline@zmax) --
                (\plinescalex*\pline@xmin,\plinescaley*\pline@ymax,\plinescalez*\pline@zmax) -- cycle;
            \path[spath/save global = xminymin#2] 
                (\plinescalex*\pline@xmin,\plinescaley*\pline@ymin,\plinescalez*\pline@zmin) -- 
                (\plinescalex*\pline@xmin,\plinescaley*\pline@ymin,\plinescalez*\pline@zmax) --
                ++(0,0,{0.25*(\plinescalez*\pline@zmax-\plinescalez*\pline@zmin)})
                ;
            \path[spath/save global = xmaxymin#2] 
                (\plinescalex*\pline@xmax,\plinescaley*\pline@ymin,\plinescalez*\pline@zmin) -- 
                (\plinescalex*\pline@xmax,\plinescaley*\pline@ymin,\plinescalez*\pline@zmax);
            \path[spath/save global = xmaxymax#2] 
                (\plinescalex*\pline@xmax,\plinescaley*\pline@ymax,\plinescalez*\pline@zmin) -- 
                (\plinescalex*\pline@xmax,\plinescaley*\pline@ymax,\plinescalez*\pline@zmax);
            \path[spath/save global = xminymax#2] 
                (\plinescalex*\pline@xmin,\plinescaley*\pline@ymax,\plinescalez*\pline@zmin) -- 
                (\plinescalex*\pline@xmin,\plinescaley*\pline@ymax,\plinescalez*\pline@zmax);
        \end{scope}
    \fi
    \endgroup
}